\section{Phase 2: oracle model comparison}
\label{sec:phase2}
This is the first direct test of the central hypothesis. Family labels are known,
data are long and low-noise, and learning imperfections are intentionally absent.

\subsection{Experiment 2A: construct M1--M5}
\paragraph{Objective.}
Construct the architectures in \cref{tab:methods} under matched data support.

\paragraph{Setup and fairness rules.}
Global models pool all clients. Clustered models use true family labels. LTI
models are fitted across the same operating data as LPV models rather than at one
favorable speed. Aggregation is uniform initially. Uncertainty weighting and
learned clusters remain outside the oracle comparison.

\paragraph{Decision gate.}
All methods must pass numerical checks and use identical train/test splits. Model
errors are decomposed by method, speed, family, client, and trajectory.

\paragraph{Results, interpretation, and decision.}
Pending.

\subsection{Experiment 2B: isolate scheduling and specialization}
\paragraph{Objective.}
Determine whether scheduling and structural specialization solve distinct and
complementary errors.

\paragraph{Predeclared expectation.}
For speed variation within one family, LPV models outperform LTI models. At fixed
speed across families, specialized models outperform global models. With both
effects present, M4 outperforms M2 and M3.

\paragraph{Setup and planned evidence.}
Three diagnostic conditions are used:
\begin{enumerate}
 \item D1: one family over the full speed envelope;
 \item D2: all families at $v_x=20\,\mathrm{m\,s^{-1}}$;
 \item D3: all families under varying speed.
\end{enumerate}
Report one-step and free-run errors, error-versus-speed curves, worst-family error,
and family gap. After development, repeat with ten independently sampled fleets.

\paragraph{Decision gate.}
M4 must improve on both M2 and M3 in D3. A useful working gate is more than 10\%
for LPV relevance under broad speed variation and approximately 5--10\% for
specialization, with practical scale and seed consistency taking precedence over
the numerical thresholds. If M3 and M4 are equivalent, clustering is unnecessary.
If M2 and M4 are equivalent, LPV scheduling is unnecessary.

\paragraph{Results, interpretation, and decision.}
Pending.

